% ============================================================
% Section 2: Background and Threat Model
% Drop this block into the main .tex file as % ============================================================
% Section 2: Background and Threat Model
% Drop this block into the main .tex file as % ============================================================
% Section 2: Background and Threat Model
% Drop this block into the main .tex file as % ============================================================
% Section 2: Background and Threat Model
% Drop this block into the main .tex file as \input{sec_threat_model}
% ============================================================

\section{Background and Threat Model}
\label{sec:threat}

\subsection{Membership Inference Attacks}

Let $D = \{(x_i, y_i)\}_{i=1}^{n}$ be a training dataset drawn i.i.d.\ from an
unknown distribution $\mathcal{P}$.
A learner trains a classifier $f_\theta : \mathcal{X} \to \Delta(\mathcal{Y})$
by minimizing empirical risk on $D$.
A membership inference attack (MIA) is an algorithm
$\mathcal{A} : \mathcal{X} \times \mathcal{Y} \times \mathcal{F} \to \{0,1\}$
that, given a target sample $(x, y)$ and query access to $f_\theta$,
outputs $1$ if $(x,y) \in D$ and $0$ otherwise.

The MIA advantage of $\mathcal{A}$ is:
\begin{equation}
  \mathrm{Adv}(\mathcal{A}, f_\theta, D)
  \;=\;
  \left|
    \Pr[\mathcal{A}(x,y,f_\theta)=1 \mid (x,y)\in D]
    -
    \Pr[\mathcal{A}(x,y,f_\theta)=1 \mid (x,y)\notin D]
  \right|.
  \label{eq:adv}
\end{equation}
We denote the dataset-level risk as the expected advantage over samples:
\begin{equation}
  \mathrm{Risk}(D) \;=\; \sup_{\mathcal{A}} \;\mathbb{E}_{(x,y)\sim\mathcal{P}}\,
  \mathrm{Adv}(\mathcal{A}, f_\theta, D),
  \label{eq:risk}
\end{equation}
where the supremum is taken over all computationally bounded adversaries.
In practice we estimate $\mathrm{Risk}(D)$ via the AUC of the best empirical
attack (Section~\ref{sec:setup}).

\subsection{Threat Model}

We consider two principals: a data holder who trains a model on $D$
and may later publish it, and an adversary who wishes to infer
membership in $D$.

\paragraph{Adversary goal.}
Given a target sample $(x,y)$ and oracle access to $f_\theta$,
determine whether $(x,y) \in D$.

\paragraph{Adversary capability.}
The adversary has black-box access to $f_\theta$: it may submit
arbitrary inputs and observe predicted class probabilities, but cannot
inspect model weights. This is the standard setting for deployed models
(APIs, MLaaS).
As a secondary experiment (Section~\ref{sec:implications}) we also consider
white-box access, and results are consistent.

\paragraph{Adversary knowledge.}
The adversary knows the data domain distribution $\mathcal{P}$
(e.g., it knows the dataset came from medical records of a particular
demographic) but does not know the exact set $D$.
This reflects realistic scenarios where domain knowledge is public but
training data is confidential.

\paragraph{Defender goal.}
Before training begins, the data holder computes a Dataset Privacy
Risk Index (DPRI) from the raw dataset $D$ using only data-distributional
statistics, with no trained model required.
DPRI provides an estimate of $\mathrm{Risk}(D)$ that can inform
(i) whether to apply differential privacy before training,
(ii) which epsilon budget is appropriate, and
(iii) whether the dataset is suitable for public model release.

\paragraph{Scope.}
Our primary analysis targets centralized training.
We discuss extensions to federated learning and DP-SGD in
Section~\ref{sec:implications}.
We do not consider attacks on the training data directly (e.g., data
poisoning). The adversary is passive.

\paragraph{Dual use.}
DPRI can also be used offensively: an adversary who has access to
candidate target datasets (e.g., a public benchmark or a leaked schema)
can compute DPRI before mounting an attack and prioritize
high-risk datasets to maximize attack efficiency.
We demonstrate this in Section~\ref{sec:adv_use} and discuss mitigations.

% ============================================================

\section{Background and Threat Model}
\label{sec:threat}

\subsection{Membership Inference Attacks}

Let $D = \{(x_i, y_i)\}_{i=1}^{n}$ be a training dataset drawn i.i.d.\ from an
unknown distribution $\mathcal{P}$.
A learner trains a classifier $f_\theta : \mathcal{X} \to \Delta(\mathcal{Y})$
by minimizing empirical risk on $D$.
A membership inference attack (MIA) is an algorithm
$\mathcal{A} : \mathcal{X} \times \mathcal{Y} \times \mathcal{F} \to \{0,1\}$
that, given a target sample $(x, y)$ and query access to $f_\theta$,
outputs $1$ if $(x,y) \in D$ and $0$ otherwise.

The MIA advantage of $\mathcal{A}$ is:
\begin{equation}
  \mathrm{Adv}(\mathcal{A}, f_\theta, D)
  \;=\;
  \left|
    \Pr[\mathcal{A}(x,y,f_\theta)=1 \mid (x,y)\in D]
    -
    \Pr[\mathcal{A}(x,y,f_\theta)=1 \mid (x,y)\notin D]
  \right|.
  \label{eq:adv}
\end{equation}
We denote the dataset-level risk as the expected advantage over samples:
\begin{equation}
  \mathrm{Risk}(D) \;=\; \sup_{\mathcal{A}} \;\mathbb{E}_{(x,y)\sim\mathcal{P}}\,
  \mathrm{Adv}(\mathcal{A}, f_\theta, D),
  \label{eq:risk}
\end{equation}
where the supremum is taken over all computationally bounded adversaries.
In practice we estimate $\mathrm{Risk}(D)$ via the AUC of the best empirical
attack (Section~\ref{sec:setup}).

\subsection{Threat Model}

We consider two principals: a data holder who trains a model on $D$
and may later publish it, and an adversary who wishes to infer
membership in $D$.

\paragraph{Adversary goal.}
Given a target sample $(x,y)$ and oracle access to $f_\theta$,
determine whether $(x,y) \in D$.

\paragraph{Adversary capability.}
The adversary has black-box access to $f_\theta$: it may submit
arbitrary inputs and observe predicted class probabilities, but cannot
inspect model weights. This is the standard setting for deployed models
(APIs, MLaaS).
As a secondary experiment (Section~\ref{sec:implications}) we also consider
white-box access, and results are consistent.

\paragraph{Adversary knowledge.}
The adversary knows the data domain distribution $\mathcal{P}$
(e.g., it knows the dataset came from medical records of a particular
demographic) but does not know the exact set $D$.
This reflects realistic scenarios where domain knowledge is public but
training data is confidential.

\paragraph{Defender goal.}
Before training begins, the data holder computes a Dataset Privacy
Risk Index (DPRI) from the raw dataset $D$ using only data-distributional
statistics, with no trained model required.
DPRI provides an estimate of $\mathrm{Risk}(D)$ that can inform
(i) whether to apply differential privacy before training,
(ii) which epsilon budget is appropriate, and
(iii) whether the dataset is suitable for public model release.

\paragraph{Scope.}
Our primary analysis targets centralized training.
We discuss extensions to federated learning and DP-SGD in
Section~\ref{sec:implications}.
We do not consider attacks on the training data directly (e.g., data
poisoning). The adversary is passive.

\paragraph{Dual use.}
DPRI can also be used offensively: an adversary who has access to
candidate target datasets (e.g., a public benchmark or a leaked schema)
can compute DPRI before mounting an attack and prioritize
high-risk datasets to maximize attack efficiency.
We demonstrate this in Section~\ref{sec:adv_use} and discuss mitigations.

% ============================================================

\section{Background and Threat Model}
\label{sec:threat}

\subsection{Membership Inference Attacks}

Let $D = \{(x_i, y_i)\}_{i=1}^{n}$ be a training dataset drawn i.i.d.\ from an
unknown distribution $\mathcal{P}$.
A learner trains a classifier $f_\theta : \mathcal{X} \to \Delta(\mathcal{Y})$
by minimizing empirical risk on $D$.
A membership inference attack (MIA) is an algorithm
$\mathcal{A} : \mathcal{X} \times \mathcal{Y} \times \mathcal{F} \to \{0,1\}$
that, given a target sample $(x, y)$ and query access to $f_\theta$,
outputs $1$ if $(x,y) \in D$ and $0$ otherwise.

The MIA advantage of $\mathcal{A}$ is:
\begin{equation}
  \mathrm{Adv}(\mathcal{A}, f_\theta, D)
  \;=\;
  \left|
    \Pr[\mathcal{A}(x,y,f_\theta)=1 \mid (x,y)\in D]
    -
    \Pr[\mathcal{A}(x,y,f_\theta)=1 \mid (x,y)\notin D]
  \right|.
  \label{eq:adv}
\end{equation}
We denote the dataset-level risk as the expected advantage over samples:
\begin{equation}
  \mathrm{Risk}(D) \;=\; \sup_{\mathcal{A}} \;\mathbb{E}_{(x,y)\sim\mathcal{P}}\,
  \mathrm{Adv}(\mathcal{A}, f_\theta, D),
  \label{eq:risk}
\end{equation}
where the supremum is taken over all computationally bounded adversaries.
In practice we estimate $\mathrm{Risk}(D)$ via the AUC of the best empirical
attack (Section~\ref{sec:setup}).

\subsection{Threat Model}

We consider two principals: a data holder who trains a model on $D$
and may later publish it, and an adversary who wishes to infer
membership in $D$.

\paragraph{Adversary goal.}
Given a target sample $(x,y)$ and oracle access to $f_\theta$,
determine whether $(x,y) \in D$.

\paragraph{Adversary capability.}
The adversary has black-box access to $f_\theta$: it may submit
arbitrary inputs and observe predicted class probabilities, but cannot
inspect model weights. This is the standard setting for deployed models
(APIs, MLaaS).
As a secondary experiment (Section~\ref{sec:implications}) we also consider
white-box access, and results are consistent.

\paragraph{Adversary knowledge.}
The adversary knows the data domain distribution $\mathcal{P}$
(e.g., it knows the dataset came from medical records of a particular
demographic) but does not know the exact set $D$.
This reflects realistic scenarios where domain knowledge is public but
training data is confidential.

\paragraph{Defender goal.}
Before training begins, the data holder computes a Dataset Privacy
Risk Index (DPRI) from the raw dataset $D$ using only data-distributional
statistics, with no trained model required.
DPRI provides an estimate of $\mathrm{Risk}(D)$ that can inform
(i) whether to apply differential privacy before training,
(ii) which epsilon budget is appropriate, and
(iii) whether the dataset is suitable for public model release.

\paragraph{Scope.}
Our primary analysis targets centralized training.
We discuss extensions to federated learning and DP-SGD in
Section~\ref{sec:implications}.
We do not consider attacks on the training data directly (e.g., data
poisoning). The adversary is passive.

\paragraph{Dual use.}
DPRI can also be used offensively: an adversary who has access to
candidate target datasets (e.g., a public benchmark or a leaked schema)
can compute DPRI before mounting an attack and prioritize
high-risk datasets to maximize attack efficiency.
We demonstrate this in Section~\ref{sec:adv_use} and discuss mitigations.

% ============================================================

\section{Background and Threat Model}
\label{sec:threat}

\subsection{Membership Inference Attacks}

Let $D = \{(x_i, y_i)\}_{i=1}^{n}$ be a training dataset drawn i.i.d.\ from an
unknown distribution $\mathcal{P}$.
A learner trains a classifier $f_\theta : \mathcal{X} \to \Delta(\mathcal{Y})$
by minimizing empirical risk on $D$.
A membership inference attack (MIA) is an algorithm
$\mathcal{A} : \mathcal{X} \times \mathcal{Y} \times \mathcal{F} \to \{0,1\}$
that, given a target sample $(x, y)$ and query access to $f_\theta$,
outputs $1$ if $(x,y) \in D$ and $0$ otherwise.

The MIA advantage of $\mathcal{A}$ is:
\begin{equation}
  \mathrm{Adv}(\mathcal{A}, f_\theta, D)
  \;=\;
  \left|
    \Pr[\mathcal{A}(x,y,f_\theta)=1 \mid (x,y)\in D]
    -
    \Pr[\mathcal{A}(x,y,f_\theta)=1 \mid (x,y)\notin D]
  \right|.
  \label{eq:adv}
\end{equation}
We denote the dataset-level risk as the expected advantage over samples:
\begin{equation}
  \mathrm{Risk}(D) \;=\; \sup_{\mathcal{A}} \;\mathbb{E}_{(x,y)\sim\mathcal{P}}\,
  \mathrm{Adv}(\mathcal{A}, f_\theta, D),
  \label{eq:risk}
\end{equation}
where the supremum is taken over all computationally bounded adversaries.
In practice we estimate $\mathrm{Risk}(D)$ via the AUC of the best empirical
attack (Section~\ref{sec:setup}).

\subsection{Threat Model}

We consider two principals: a data holder who trains a model on $D$
and may later publish it, and an adversary who wishes to infer
membership in $D$.

\paragraph{Adversary goal.}
Given a target sample $(x,y)$ and oracle access to $f_\theta$,
determine whether $(x,y) \in D$.

\paragraph{Adversary capability.}
The adversary has black-box access to $f_\theta$: it may submit
arbitrary inputs and observe predicted class probabilities, but cannot
inspect model weights. This is the standard setting for deployed models
(APIs, MLaaS).
As a secondary experiment (Section~\ref{sec:implications}) we also consider
white-box access, and results are consistent.

\paragraph{Adversary knowledge.}
The adversary knows the data domain distribution $\mathcal{P}$
(e.g., it knows the dataset came from medical records of a particular
demographic) but does not know the exact set $D$.
This reflects realistic scenarios where domain knowledge is public but
training data is confidential.

\paragraph{Defender goal.}
Before training begins, the data holder computes a Dataset Privacy
Risk Index (DPRI) from the raw dataset $D$ using only data-distributional
statistics, with no trained model required.
DPRI provides an estimate of $\mathrm{Risk}(D)$ that can inform
(i) whether to apply differential privacy before training,
(ii) which epsilon budget is appropriate, and
(iii) whether the dataset is suitable for public model release.

\paragraph{Scope.}
Our primary analysis targets centralized training.
We discuss extensions to federated learning and DP-SGD in
Section~\ref{sec:implications}.
We do not consider attacks on the training data directly (e.g., data
poisoning). The adversary is passive.

\paragraph{Dual use.}
DPRI can also be used offensively: an adversary who has access to
candidate target datasets (e.g., a public benchmark or a leaked schema)
can compute DPRI before mounting an attack and prioritize
high-risk datasets to maximize attack efficiency.
We demonstrate this in Section~\ref{sec:adv_use} and discuss mitigations.
